\section{Deterministic row occupancy and the long-lcm frontier}
\label{sec:tpc21-occupancy-frontier}

This section uses one feature which is absent from an arbitrary fiber
vector: the fibers come from rows \(m=\ell d\), with
\(\ell\asymp L\), \(d\asymp D\), and \(L>2R\).  It gives a genuine
anti-concentration estimate, but the resulting discrepancy bound reaches
only the already-completable side of the joint-period boundary.  The
Ramanujan mean is handled separately and more strongly in
\cref{sec:tpc21-mean-closure}.

\subsection{Safe restoration of the Cartesian row packet}

Put
\begin{equation}
 Z_m^{\mathrm{sh}}(j)
 =\sum_{\substack{u\le R\\(u,mH)=1}}
 b_R(u)\left(\one_{u\mid mj+h}-\frac1u\right).
 \label{eq:tpc21-short-signal}
\end{equation}
Only the short coefficient \(b_R(u)\one_{u\le R}\) is used here.

\begin{lemma}[Pointwise size of the isolated short signal]
\label{lem:tpc21-short-pointwise}
Uniformly on the TPC row and orbit support, for every fixed
\(\epsilon>0\),
\begin{equation}
 |Z_m^{\mathrm{sh}}(j)|\ll_{\epsilon,h}X^\epsilon.
 \label{eq:tpc21-short-pointwise}
\end{equation}
\end{lemma}

\begin{proof}
The coefficient is squarefree supported and
\begin{equation}
 \max_{u\le R}|b_R(u)|\ll (\log(2R))^2,
 \qquad
 \sum_{u\le R}\frac{|b_R(u)|}{u}\ll(\log(2R))^2.
 \label{eq:tpc21-b-short-bounds}
\end{equation}
Indeed, the first estimate follows from the positive formula for
\(\lambda'_R\), while the second follows from
\(b_R(u)=-\mu(u)\log u-\lambda'_R(u)\) and the Euler-product estimate
\(\sum_{u\le R}|\lambda'_R(u)|/u\ll(\log(2R))^2\).
Consequently
\[
 |Z_m^{\mathrm{sh}}(j)|
 \le \tau(mj+h)\max_{u\le R}|b_R(u)|
     +\sum_{u\le R}\frac{|b_R(u)|}{u}.
\]
The targets are \(O(X)\), so the standard divisor bound proves
\eqref{eq:tpc21-short-pointwise}.
\end{proof}

\begin{proposition}[Scalar-safe unpruning]
\label{prop:tpc21-scalar-unpruning}
Fix \(\kappa,\lambda_0>0\), assume \(L\ge X^{\lambda_0}\), and let the
row weights be polylogarithmically bounded.  In a scalar short--short
correlation, the diagonal, same-source, near-determinant, and large
row-gcd pairs may be restored at total cost
\begin{align}
 O_\epsilon\!\left(
 X^{1+\epsilon}
 +XQ\{L^{-1}+X^{-\kappa}\}X^\epsilon
 \right).
 \label{eq:tpc21-unpruning-error}
\end{align}
In particular, after taking \(\epsilon>0\) sufficiently small, this is
\(O_A(XQ(\log X)^{-A})\) for every fixed \(A>0\).
\end{proposition}

\begin{proof}
Apply \cref{lem:tpc21-short-pointwise} to both short signals and repeat
the tuple counts in the TPC-18 degenerate-channel removal
\citep{WangTPC18}.  The diagonal
has \(O(LDJ)=O(X)\) tuples.  The same-source channel has
\(O(LD^2J)=O(XQ/L)\) tuples, the near channel has
\(O(XQ(X^{-\kappa}+L^{-1}))\), and the large-gcd channel has
\(O(XQX^{-\kappa})\).  Source logarithms and bounded smooth weights are
absorbed into \(X^\epsilon\).  Since \(Q\ge L\ge X^{\lambda_0}\), the
diagonal also has a fixed power saving relative to \(XQ\).
\end{proof}

\begin{remark}[The order of operations matters]
\label{rem:tpc21-unpruning-caveat}
\Cref{prop:tpc21-scalar-unpruning} permits restoration of the full
Cartesian row set in the \emph{total scalar short--short packet}.  It
does not say that every individual \(q\)-fiber, Poisson frequency, or
mean/discrepancy piece of the excluded rows is small.  Thus one must
unprune before the fiber decomposition and keep
\eqref{eq:tpc21-unpruning-error} outside it; coefficientwise unpruning
would be an unjustified stronger assertion.
\end{remark}

\subsection{A row-fiber occupancy theorem}

Let \(\mathscr A\) be the full row family
\begin{equation}
 \alpha=(\ell,d),\qquad
 \ell\in[L/2,2L]\text{ prime},\qquad
 d\in[D,2D]\cap[1,V],\qquad m_\alpha=\ell d,
 \label{eq:tpc21-row-family}
\end{equation}
with any rowwise coprimality masks inserted into the coefficients.  A
smooth separation of the original two-row weight gives bilinear packets
\(x_\alpha y_\beta F(j/J)\).  We assume
\begin{equation}
 |x_\alpha|\le W_x,\qquad |y_\alpha|\le W_y,
 \label{eq:tpc21-row-weight-size}
\end{equation}
where \(W_x,W_y\) are polylogarithmic in the application.  For
squarefree \(u\le R\), \((u,H)=1\), define
\begin{equation}
 A_u^x(\kappa)
 =\sum_{\substack{\alpha\in\mathscr A\\
 m_\alpha\kappa\equiv-h\ (\bmod u)}}x_\alpha,
 \qquad \kappa\in\mathcal U_u,
 \label{eq:tpc21-one-row-fiber}
\end{equation}
and define \(A_u^y\) similarly.

\begin{lemma}[Deterministic row anti-concentration]
\label{lem:tpc21-row-occupancy}
Uniformly for \(u\le R\),
\begin{align}
 \|A_u^x\|_\infty
 &\ll W_x\frac Qu,
 \label{eq:tpc21-row-linfty}\\
 \sum_{\kappa\in\mathcal U_u}|A_u^x(\kappa)|
 &\ll W_xQ,
 \label{eq:tpc21-row-lone}\\
 \|A_u^x\|_2
 &\ll W_x\frac Q{\sqrt u}.
 \label{eq:tpc21-row-ltwo}
\end{align}
The same estimates hold with \(x\) replaced by \(y\).
\end{lemma}

\begin{proof}
Fix \(d\) and \(\kappa\).  The congruence in
\eqref{eq:tpc21-one-row-fiber} confines \(\ell\) to one class modulo
\(u\).  An interval of length \(O(L)\) contains
\(O(L/u+1)=O(L/u)\) integers in that class because \(L>2R\ge2u\).
There are \(O(D)\) choices of \(d\), proving
\eqref{eq:tpc21-row-linfty}.  Every admissible row belongs to exactly
one class \(\kappa\), so summing absolute values before grouping proves
\eqref{eq:tpc21-row-lone}.  Finally,
\(\|A\|_2^2\le\|A\|_\infty\|A\|_1\), which gives
\eqref{eq:tpc21-row-ltwo}.
\end{proof}

Let \(\vartheta\) be supported on squarefree \(u\le R\), and suppose
\begin{equation}
 |\vartheta(u)|\le\Theta_R.
 \label{eq:tpc21-theta-size}
\end{equation}
Both \(\lambda'_R(u)\) and \(b_R(u)\one_{u\le R}\) satisfy this with
\(\Theta_R\ll(\log(2R))^2\).  For squarefree \(q\) with \((q,H)=1\),
define the full
bilinear lcm fiber
\begin{equation}
 Z_q^{x,y}(\kappa)
 =\sum_{\substack{u,v\le R\\{}[u,v]=q}}
 \vartheta(u)\vartheta(v)
 A_u^x(\kappa\bmod u)A_v^y(\kappa\bmod v),
 \qquad \kappa\in\mathcal U_q.
 \label{eq:tpc21-full-lcm-fiber}
\end{equation}
No complex conjugation is implicit in \eqref{eq:tpc21-full-lcm-fiber};
after Mellin separation the packet is bilinear.  A Hermitian square is
the special case in which the second factor is conjugated.

\begin{theorem}[Actual-fiber sup norm and variance]
\label{thm:tpc21-occupancy-variance}
Put
\begin{equation}
 \mathfrak g(q)=\prod_{p\mid q}\left(2+\frac1p\right).
 \label{eq:tpc21-occupancy-gq}
\end{equation}
Then
\begin{align}
 \|Z_q^{x,y}\|_\infty
 &\ll W_xW_y\Theta_R^2Q^2\frac{\mathfrak g(q)}q,
 \label{eq:tpc21-fiber-linfty}\\
 |\overline Z_q^{x,y}|
 &\ll W_xW_y\Theta_R^2Q^2\frac{\mathfrak g(q)}q,
 \label{eq:tpc21-fiber-mean-bound}\\
 \|D_q^{x,y}\|_2
 &\ll W_xW_y\Theta_R^2Q^2
       \frac{\mathfrak g(q)}{\sqrt q},
 \label{eq:tpc21-fiber-discrepancy-bound}
\end{align}
where \(D_q^{x,y}=Z_q^{x,y}-\overline Z_q^{x,y}\) and the mean is over
\(\mathcal U_q\).  The same three estimates hold after multiplying
each row pair in \eqref{eq:tpc21-full-lcm-fiber} by any Schur mask of
modulus at most \(1\), in particular by the generic mask
\eqref{eq:tpc21-generic-mask}.
\end{theorem}

\begin{proof}
For squarefree \(q\), independent assignment of every prime divisor to
\(u\) only, \(v\) only, or both gives the exact identity
\begin{equation}
 \sum_{\substack{u,v\mid q\\{}[u,v]=q}}\frac1{uv}
 =\prod_{p\mid q}\left(\frac2p+\frac1{p^2}\right)
 =\frac1q\prod_{p\mid q}\left(2+\frac1p\right).
 \label{eq:tpc21-exact-lcm-product}
\end{equation}
The restrictions \(u,v\le R\) can only decrease the positive sum.
Insert \eqref{eq:tpc21-row-linfty} in
\eqref{eq:tpc21-full-lcm-fiber} and use
\eqref{eq:tpc21-exact-lcm-product} to obtain
\eqref{eq:tpc21-fiber-linfty}.  If a Schur mask
\(\mathfrak m(\alpha,\beta)\) with \(|\mathfrak m|\le1\) is inserted,
the absolute value of its \((u,v,\kappa)\)-fiber is at most
\[
 \left(\sum_{\alpha\ {\rm in\ the}\ u{\rm\mbox{-}fiber}}|x_\alpha|\right)
 \left(\sum_{\beta\ {\rm in\ the}\ v{\rm\mbox{-}fiber}}|y_\beta|\right),
\]
so exactly the same majorant applies.  The mean bound is immediate.
Finally,
subtraction of the mean is orthogonal in
\(\ell^2(\mathcal U_q)\), and therefore
\[
 \|D_q^{x,y}\|_2\le\|Z_q^{x,y}\|_2
 \le\sqrt{\varphi(q)}\,\|Z_q^{x,y}\|_\infty.
\]
This proves \eqref{eq:tpc21-fiber-discrepancy-bound}.
\end{proof}

\subsection{What the occupancy bound buys}

Let \(q\asymp Y\) be a dyadic denominator block and retain the
low-frequency window of TPC-20.  Assume \(J\ge4HX^{\epsilon_0}\), so
that the positive and negative retained frequencies do not wrap modulo
\(q\), and assume the polynomial fiber-mass condition needed to remove
the rapidly decaying Fourier tail.

\begin{corollary}[Dyadic discrepancy consequence]
\label{cor:tpc21-dyadic-occupancy}
For fixed \(H\), with a summed Fourier tail smaller than every power of
\(X\),
\begin{equation}
 \sum_{q\asymp Y}|\mathcal L_q^{\mathrm{disc}}|
 \ll_{F,H,\epsilon_0}
 X^{\epsilon_0/2}W_xW_y\Theta_R^2
 Q^2\sqrt{JY}\log(2Y),
 \label{eq:tpc21-dyadic-discrepancy}
\end{equation}
Relative to the
natural scale \(XQ=Q^2J\), the right-hand side is
\begin{equation}
 \ll X^{\epsilon_0/2}W_xW_y\Theta_R^2\log(2Y)
 \sqrt{\frac YJ}.
 \label{eq:tpc21-occupancy-ratio}
\end{equation}
Consequently this method gives a fixed power saving when
\(Y\le JX^{-\eta}\), after taking \(0<\epsilon_0<\eta\), but gives no
fixed power saving in a genuinely long block \(Y\ge JX^\eta\).
\end{corollary}

\begin{proof}
The restricted-frequency Parseval estimate of TPC-20 gives
\[
 |\mathcal L_q^{\mathrm{disc}}|
 \ll X^{\epsilon_0/2}\sqrt J\,\|D_q^{x,y}\|_2.
\]
Moreover, for \(\Re s>1\),
\[
 \sum_{q\ge1}\frac{\mu^2(q)\mathfrak g(q)}{q^s}
 =\zeta(s)^2\mathcal H(s),
\]
where the local factor of \(\mathcal H\) is
\[
 \left(1+\left(2+\frac1p\right)p^{-s}\right)(1-p^{-s})^2.
\]
After the linear terms \(2p^{-s}\) cancel, the remaining Euler product
is absolutely convergent for \(\Re s>1/2\).  If
\(\mathcal H(s)=\sum_n h(n)n^{-s}\), then in particular
\(\sum_n|h(n)|/n<\infty\).  Dirichlet convolution with the divisor
function, followed by
\(\sum_{n\le T}\tau(n)\ll T\log(2T)\), gives
\begin{equation}
 \sum_{q\le T}\mu^2(q)\mathfrak g(q)\ll T\log(2T),
 \qquad
 \sum_{q\asymp Y}\frac{\mu^2(q)\mathfrak g(q)}{\sqrt q}
 \ll\sqrt Y\log(2Y).
 \label{eq:tpc21-gq-average}
\end{equation}
Sum \eqref{eq:tpc21-fiber-discrepancy-bound} and use
\eqref{eq:tpc21-gq-average}.  Division by \(Q^2J\) proves
\eqref{eq:tpc21-occupancy-ratio}.
\end{proof}

\begin{remark}[A normalization gain, not a long-lcm estimate]
\label{rem:tpc21-occupancy-frontier}
The factor \(q^{-1/2}\) in
\eqref{eq:tpc21-fiber-discrepancy-bound} is a genuine improvement over
treating \(Z_q(\kappa)\) as an arbitrary fiber vector.  Nevertheless,
\eqref{eq:tpc21-occupancy-ratio} crosses the natural scale precisely at
\(Y\asymp J\), which is also the joint-period completion boundary up to
fixed powers of \(X\) and the fixed factor \(H\).  Thus deterministic
row occupancy does not penetrate the long-lcm gate.  This is a limit of
the stated estimate, not a lower bound for the arithmetic packet.
\end{remark}

The mean term behaves better: \cref{thm:tpc21-summed-mean-closure}
closes it absolutely for the original masked packet, without restoring
the Cartesian row family or invoking the occupancy estimate.
